\section{Discussion}

The results indicate that adaptive allocation can improve single-photon LiDAR sensing efficiency without modifying the underlying transmitter or detector hardware. The proposed method differs from purely posterior-based selection because it incorporates a predicted signal-to-background ratio derived from the optical forward model. Regions with high uncertainty but extremely weak expected optical contrast are therefore not always prioritized.

The largest practical benefit is the reduction in photons and sensing steps required to obtain a confident target decision. This is relevant to photon-limited reconnaissance systems in which laser energy, dwell time, detector saturation, or platform power must be constrained.

The study remains a numerical short communication rather than a deployment demonstration. All experiments were simulation based, and the atmospheric and target models use simplified homogeneous assumptions. The evaluation included one target in every episode; consequently, the reported zero false-focus rate is not equivalent to a dedicated target-absent false-alarm probability.

The information-gain score is also a physics-guided heuristic rather than an exact mutual-information optimization. Nevertheless, the compact formulation is computationally inexpensive, interpretable, and directly linked to measurable optical quantities.

Future work should evaluate target-absent scenes, nonhomogeneous atmospheric channels, moving targets, multiple targets, and experimental single-photon LiDAR measurements.
